\labsection{4}{Regression and time-aware validation}{sec:lab04-regression}

\begin{labproject}{Linear models under ordinary and chronological validation}
\textbf{Coverage.} Chapter 9.\\
\textbf{Core data.} Medical Cost Personal Data. Bike Sharing is a transfer activity within the same lab.\\
\textbf{Goal.} Compare ordinary linear regression with ridge regularization, inspect residual behavior, and contrast random splitting with a chronological holdout when time ordering matters.

\begin{labmode}
The main learning objective is not a particular regression score. It is choosing a validation design that matches how future predictions will actually be made.
\end{labmode}

\medskip\noindent\textbf{Part A --- Linear and ridge regression.}\par
\begin{lstlisting}[style=mlpython]
from pathlib import Path
import numpy as np
import pandas as pd
import matplotlib.pyplot as plt
from sklearn.compose import ColumnTransformer
from sklearn.impute import SimpleImputer
from sklearn.linear_model import LinearRegression, Ridge
from sklearn.metrics import mean_absolute_error, mean_squared_error, r2_score
from sklearn.model_selection import train_test_split
from sklearn.pipeline import Pipeline
from sklearn.preprocessing import OneHotEncoder, StandardScaler
\end{lstlisting}

\begin{lstlisting}[style=mlpython]
PROJECT_ROOT = Path.cwd()
train_path = PROJECT_ROOT / "all_datasets" / "medical_cost_personal_dataset" / "insurance.csv"
df = pd.read_csv(train_path)
\end{lstlisting}

\begin{lstlisting}[style=mlpython]
X = df.drop(columns="charges")
y = df["charges"]
num_cols = ["age", "bmi", "children"]
cat_cols = ["sex", "smoker", "region"]

preprocess = ColumnTransformer([
    ("num", Pipeline([
        ("imputer", SimpleImputer(strategy="median")),
        ("scale", StandardScaler()),
    ]), num_cols),
    ("cat", Pipeline([
        ("imputer", SimpleImputer(strategy="most_frequent")),
        ("onehot", OneHotEncoder(handle_unknown="ignore")),
    ]), cat_cols),
])

X_train, X_valid, y_train, y_valid = train_test_split(
    X, y, test_size=0.20, random_state=42
)

predictions = {}
\end{lstlisting}

\begin{lstlisting}[style=mlpython]
for name, estimator in {
    "ordinary least squares": LinearRegression(),
    "ridge": Ridge(alpha=10.0),
}.items():
    model = Pipeline([("preprocess", preprocess), ("model", estimator)])
    model.fit(X_train, y_train)
    pred = model.predict(X_valid)
    predictions[name] = pred
    print("\n", name)
    print("MAE:", round(mean_absolute_error(y_valid, pred), 2))
    print("RMSE:", round(np.sqrt(mean_squared_error(y_valid, pred)), 2))
    print("R2:", round(r2_score(y_valid, pred), 3))
\end{lstlisting}

\begin{lstlisting}[style=mlpython]
# Visual diagnostics for the regularized model.
ridge_pred = predictions["ridge"]
plt.figure(figsize=(6, 4))
plt.scatter(y_valid, ridge_pred, alpha=0.65)
lims = [min(y_valid.min(), ridge_pred.min()), max(y_valid.max(), ridge_pred.max())]
plt.plot(lims, lims, linestyle="--")
plt.xlabel("Actual charges")
plt.ylabel("Predicted charges")
plt.title("Ridge validation: actual versus predicted")
plt.tight_layout()
plt.show()
\end{lstlisting}

\begin{lstlisting}[style=mlpython]
residual = y_valid.to_numpy() - ridge_pred
plt.figure(figsize=(6, 4))
plt.scatter(ridge_pred, residual, alpha=0.65)
plt.axhline(0, linestyle="--")
plt.xlabel("Predicted charges")
plt.ylabel("Residual")
plt.title("Ridge validation residual plot")
plt.tight_layout()
plt.show()
\end{lstlisting}

\medskip\noindent\textbf{Part B --- Optional chronological-validation transfer.}\par
\begin{lstlisting}[style=mlpython]
from pathlib import Path
import matplotlib.pyplot as plt
import numpy as np
import pandas as pd
from sklearn.compose import ColumnTransformer
from sklearn.ensemble import RandomForestRegressor
from sklearn.metrics import mean_absolute_error, mean_squared_error
from sklearn.pipeline import Pipeline
from sklearn.preprocessing import OneHotEncoder
\end{lstlisting}

\begin{lstlisting}[style=mlpython]
PROJECT_ROOT = Path.cwd()
train_path = PROJECT_ROOT / "all_datasets" / "bike_sharing_dataset" / "day.csv"
df = pd.read_csv(train_path)
df["dteday"] = pd.to_datetime(df["dteday"])
df = df.sort_values("dteday").reset_index(drop=True)
\end{lstlisting}

\begin{lstlisting}[style=mlpython]
# cnt = casual + registered, so casual and registered would leak the label.
features = [
    "season", "yr", "mnth", "holiday", "weekday", "workingday",
    "weathersit", "temp", "atemp", "hum", "windspeed",
]
X = df[features]
y = df["cnt"]
\end{lstlisting}

\begin{lstlisting}[style=mlpython]
# Chronological development holdout: earlier dates train, later dates validate.
cut = int(len(df) * 0.80)
X_train, X_valid = X.iloc[:cut], X.iloc[cut:]
y_train, y_valid = y.iloc[:cut], y.iloc[cut:]
dates_valid = df["dteday"].iloc[cut:]

cat_cols = ["season", "yr", "mnth", "holiday", "weekday", "workingday", "weathersit"]
num_cols = ["temp", "atemp", "hum", "windspeed"]
preprocess = ColumnTransformer([
    ("cat", OneHotEncoder(handle_unknown="ignore"), cat_cols),
    ("num", "passthrough", num_cols),
])
model = Pipeline([
    ("preprocess", preprocess),
    ("model", RandomForestRegressor(
        n_estimators=400, min_samples_leaf=2, random_state=42, n_jobs=-1
    )),
])
model.fit(X_train, y_train)
pred = model.predict(X_valid)

print("MAE:", round(mean_absolute_error(y_valid, pred), 1))
print("RMSE:", round(np.sqrt(mean_squared_error(y_valid, pred)), 1))
\end{lstlisting}

\begin{lstlisting}[style=mlpython]
plt.figure(figsize=(8, 4))
plt.plot(dates_valid, y_valid.to_numpy(), label="actual")
plt.plot(dates_valid, pred, label="predicted")
plt.xlabel("Date")
plt.ylabel("Daily rentals")
plt.title("Chronological validation: actual versus predicted demand")
plt.legend()
plt.tight_layout()
plt.show()
\end{lstlisting}


\end{labproject}

\begin{labdeliverable}
Submit linear-versus-ridge validation metrics, one residual diagnostic, and a short comparison between random and chronological validation on the bike-sharing transfer activity. Explain which split better matches deployment and why.
\end{labdeliverable}
